\CQUsection{相关理论基础}

\CQUsubsection{医学图像分割基本理论}
医学图像分割的目标是在像素或体素层面对目标解剖结构进行精确定位，其本质是一个密集预测问题。与自然图像中的目标检测相比，CTA 体数据具有三维空间连续、前景占比极低、病灶形态差异显著等特点，因此模型不仅需要学习局部边缘与纹理，还要保持跨层空间一致性。对于颅内动脉瘤任务而言，分割结果既可作为病灶可视化依据，也可进一步支持体积、长径、形态不规则度等衍生指标的计算，从而成为后续风险评估的重要输入。

在经典监督学习框架下，训练样本由输入图像 $X$ 与金标准标签 $Y$ 构成，模型学习从 $X$ 到 $Y$ 的映射函数 $f_{\theta}(X)$。由于医学图像标注成本高、类别分布不平衡、成像条件复杂，实际训练中往往需要在网络结构、损失函数与采样机制之间协同设计，以保证模型既具有表达能力，又能稳定收敛。

\CQUsubsection{卷积神经网络与 U-Net 结构}
卷积神经网络（Convolutional Neural Network, CNN）通过局部感受野、权值共享和层级特征提取机制，在图像分析任务中取得了广泛应用。对于分割问题，FCN（Fully Convolutional Network）首先将分类网络推广到像素级预测，使网络能够直接生成密集输出图\cite{long2015}；而 U-Net 在此基础上引入对称的编码器—解码器结构及跳跃连接，使浅层细节信息与深层语义信息得以融合，显著提升了医学图像分割精度\cite{unet}。

三维医学图像具有明显的层间连续性，因此三维卷积网络能够更自然地利用上下文信息。3D U-Net 将二维卷积替换为三维卷积与三维上采样，可同时建模切片内与切片间的空间关系\cite{unet3d}。然而，三维卷积也会带来更高的显存与计算开销，因此在实际工程实现中通常需要结合补丁采样、流式加载和混合精度等策略来控制资源消耗。

\CQUsubsection{注意力机制与 Attention U-Net}
注意力机制的核心思想是让模型在大量候选特征中自适应地突出与当前任务更相关的部分。对于颅内动脉瘤分割任务，病灶体积小、背景血管复杂，若直接拼接所有跳跃连接特征，容易将大量与病灶无关的响应传递至解码器，增加误检风险。Attention U-Net 在跳跃连接处加入门控模块，对编码器特征进行筛选，使解码器更关注目标区域\cite{attention_unet}。

从形式上看，注意力门控通过当前解码特征与同尺度编码特征的联合映射生成一张权重图，再对编码特征进行加权。该机制可以理解为一种空间相关性筛选：当某位置更可能属于动脉瘤或其邻近血管区域时，权重会相对提高；当某位置更可能是背景时，权重会被抑制。这样既有助于减少无关背景干扰，又保留了 U-Net 对局部边界细节的表达优势。

\CQUsubsection{类别不平衡与损失函数基础}
在颅内动脉瘤任务中，目标体素通常只占整例 CTA 体积中的极小比例，属于典型的类别极不平衡问题。若仅使用标准交叉熵损失，模型很容易学到“全部预测为背景”这一平凡解，从而在总体准确率上表现较高，却难以真正检出病灶。因此，需要引入更关注前景重叠程度的区域型损失。

Dice 系数是衡量预测区域与真实区域重叠程度的常用指标，其对应的 Dice 损失能够直接优化前景区域的一致性。对于预测概率 $p_i$ 与标签 $g_i$，Dice 损失可写为：
\begin{equation}
\mathcal{L}_{Dice}=1-\frac{2\sum_i p_i g_i+\epsilon}{\sum_i p_i+\sum_i g_i+\epsilon},
\end{equation}
其中 $\epsilon$ 为平滑项。交叉熵损失则从逐体素分类角度约束模型输出，强调概率预测的数值稳定性。将二者组合能够兼顾局部分类准确性与整体区域重叠度，这也是本文分割模块损失设计的理论基础。

对于风险预测任务，破裂与未破裂样本往往同样存在比例失衡，因此二分类训练中需要通过正类加权、阈值搜索等方式提高模型对少数类的敏感性。相比固定阈值 0.5，基于验证集的阈值搜索能够在精确率与召回率之间取得更符合任务目标的平衡。

\CQUsubsection{表格学习与特征表示理论}
除影像体数据外，临床信息、人口统计学特征与既往病史通常以结构化表格形式存在。传统表格学习方法如逻辑回归、支持向量机和随机森林在变量规模有限时具有一定优势，但在数值特征与类别特征混合、非线性交互显著的场景下，模型表达能力可能受到限制。

近年来，深度表格学习开始借鉴自然语言处理中的 token 化思想：将每个数值字段、类别字段都映射为统一维度的隐向量，再利用注意力机制建模字段之间的交互关系。其优点在于：一方面避免了人工构造高阶交叉项，另一方面能够统一处理连续变量与离散变量，为多模态融合提供了更灵活的表示空间。本文风险模型即基于这一思想，将数值 token 与类别 token 送入交叉注意力模块，以学习临床变量之间的潜在关联。

\CQUsubsection{评价指标基础}
分割与分类任务的评价重点不同。分割侧常用 Dice 系数、敏感度、特异度、Hausdorff 距离等指标，其中 Dice 强调重叠区域的一致性，适合评价前景区域的整体拟合程度。分类与风险预测侧则常用准确率、精确率、召回率、F1 值与 AUC。AUC 反映模型对正负样本排序能力的整体水平，而 F1 值综合考虑精确率与召回率，适合在类别不平衡场景下衡量模型的综合表现。

需要指出的是，对于极度稀疏的医学分割任务，单独依赖总体准确率往往会产生误导，因为背景体素数量远大于前景体素数量。正因如此，本文在实验分析中同时报告 Dice、体素级 F1、多阈值统计与基于直方图近似的 voxel AUC，以尽可能全面地反映模型性能。

\CQUsubsection{本章小结}
本章围绕本文后续方法设计所依赖的理论基础进行了概述，主要包括三维医学图像分割的基本问题、CNN 与 U-Net 结构、注意力门控机制、类别不平衡下的损失函数设计、表格特征表示与交叉注意力思想，以及分割和风险预测常用评价指标。上述理论共同构成了本文系统实现与实验分析的出发点，为下一章的模型设计与工程实现提供了方法依据。
